%% kapa: exam=1 round=32 year=2021 period=2 form=A subject=appraisal_law subject_ko=감정평가관계법규 seq=35 number=35
%% source: Q-Net 감정평가사 기출문제 (cst00309) artlSeq=5208397 file=extracted/1008/5208397/A형/2021년도 제32회 감정평가사 1차 2교시 A형.pdf
%% conversion: kordoc 4.24.51 -> tools/gen1.py (auto); review=proofread-sonnet; flags=none
\begin{question}{35}
부동산등기법령상 가등기에 관한 설명으로 옳지 않은 것은?
\begin{choices}
  \item 가등기를 명하는 가처분명령의 관할법원은 부동산의 소재지를 관할하는 지방법원이다.
  \item 가등기권리자는 가등기를 명하는 법원의 가처분명령이 있을 때에는 단독으로 가등기를 신청할 수 있다.
  \item 가등기를 명하는 가처분명령의 신청을 각하하는 결정에 대하여는 즉시항고를 할 수 있다.
  \item 가등기에 의한 본등기를 한 경우 본등기의 순위는 가등기의 순위에 따른다.
  \item 가등기의무자는 가등기명의인의 동의 없이도 단독으로 가등기의 말소를 신청할 수 있다.
\end{choices}
\end{question}
\answer{5}
